% Section 3: Background and Method (~1.5 pages, ~1500 words)
% Write FIRST. This defines the technical contribution; everything else flows from it.

\section{Background and Method}
\label{sec:method}

\subsection{Preliminaries}
\label{sec:prelim}

% Define MDP/POMDP notation, Watkins Q(lambda), Pokemon Showdown as POMDP,
% hierarchical action decomposition (20-tuple master state, 17-tuple sub-state),
% heuristic prior h(s, a).

We model a Pokemon battle as a partially observable Markov decision process (POMDP)
$(\mathcal{S}, \mathcal{A}, P, r, \gamma, \Omega)$ where $\mathcal{S}$ is the set of
joint game states, $\mathcal{A}$ is the player's action set (move 1 through 4, plus a
switch action), $P$ is a stochastic transition function determined by the Showdown
simulator, $r$ is the dense reward (HP delta, faint delta, terminal $\pm 1$),
$\gamma$ is the discount factor, and $\Omega : \mathcal{S} \to \mathcal{O}$ is the
observation function that hides opponent team information.

We train each member with Watkins $Q(\lambda)$ \citep{watkins1989, sutton2018}:
\begin{align}
  \delta_t       &= r_{t+1} + \gamma \max_{a'} Q(s_{t+1}, a') - Q(s_t, a_t) \\
  e_t(s, a)      &= \begin{cases} 1, & (s, a) = (s_t, a_t) \\
                                  \gamma \lambda \, e_{t-1}(s, a), & a_t = \arg\max_{a'} Q(s_t, a') \\
                                  0, & \text{otherwise} \end{cases} \\
  Q(s, a)        &\leftarrow Q(s, a) + \alpha \delta_t \, e_t(s, a)
\end{align}
with replacing traces and the trace cleared after non-greedy exploration.

% Hierarchical decomposition (from V5)
We adopt the hierarchical action decomposition introduced in our prior work
\citep{singh2026v5}: a \emph{master} agent observes a 20-tuple state and chooses
between attacking (one of four moves) and switching; if it chooses to switch,
a \emph{sub-agent} observes a 17-tuple state and selects which Pokemon to switch to.
Both agents are trained jointly via the shared reward signal. We use the master's
Q-table $Q^{\text{m}}$ and the sub-agent's Q-table $Q^{\text{s}}$ throughout.

% Heuristic prior h(s, a)
For unseen state-action pairs at inference time, we use the same domain heuristic
$h(s, a)$ that V5 used to ``smart-initialize'' the Q-tables: it is a deterministic
function of the battle state, combining type effectiveness, STAB, move power, and
accuracy. Crucially, $h(s, a)$ depends only on $s$ and the action's type signature,
not on any trained Q-table.

\subsection{The ensemble inference rule}
\label{sec:method_inference}

% This is the technical contribution. State formally.
% Define: Q-matrix at decision time, three combination rules, read-only property.

Given $K$ trained Q-tables $\{Q^k_{\text{m}}, Q^k_{\text{s}}\}_{k=1}^K$, we choose
the agent's action at state $s$ as follows.

\paragraph{Action enumeration.}
Following the master/sub-agent decomposition, we first enumerate the possible
\emph{master} actions $\mathcal{A}_{\text{m}}(s)$ available at state $s$
(legal moves of the active Pokemon, plus a switch action when force-switch is not
required). We then construct the \emph{Q-matrix} $\mathbf{Q} \in \mathbb{R}^{K \times |\mathcal{A}_{\text{m}}|}$
whose entry $\mathbf{Q}_{k, a}$ is the master Q-value of action $a$ at state $s$ under
member $k$, with heuristic fallback for unseen entries:
\begin{equation}
  \mathbf{Q}_{k, a} := \begin{cases}
    Q^k_{\text{m}}(s, a), & \text{if } (s, a) \in \text{dom}(Q^k_{\text{m}}) \\
    h(s, a),              & \text{otherwise.}
  \end{cases}
  \label{eq:q_matrix}
\end{equation}

\paragraph{Combination rules.}
We compare three rules for combining the $K$ rows of $\mathbf{Q}$ into a single
chosen action $a^* \in \mathcal{A}_{\text{m}}$.

\textbf{Soft voting} averages Q-values column-wise and picks the argmax:
\begin{equation}
  a^*_{\text{soft}} = \arg\max_{a} \frac{1}{K} \sum_{k=1}^K \mathbf{Q}_{k, a}.
\end{equation}

\textbf{Hard voting} (plurality) lets each member vote for its preferred action:
\begin{equation}
  a^*_{\text{hard}} = \arg\max_{a} \big|\{ k : a = \arg\max_{a'} \mathbf{Q}_{k, a'}\}\big|,
\end{equation}
with ties broken by mean Q across members.

\textbf{Confidence-weighted voting} assigns each member a weight equal to the gap
between its top Q-value and its row mean:
\begin{align}
  w_k &:= \max\!\left( \max_a \mathbf{Q}_{k, a} - \tfrac{1}{|\mathcal{A}_{\text{m}}|} \sum_a \mathbf{Q}_{k, a},\ 0 \right), \\
  a^*_{\text{conf}} &:= \arg\max_{a} \frac{\sum_k w_k \mathbf{Q}_{k, a}}{\sum_k w_k}.
\end{align}
When $\sum_k w_k = 0$, we fall back to soft voting.

\paragraph{Recursion into the sub-agent.}
If the chosen master action is \emph{switch}, we repeat the same procedure on the
sub-agent Q-tables $\{Q^k_{\text{s}}\}_{k=1}^K$ over the available switch candidates.

% Key property
\paragraph{Read-only property.}
None of the three rules modifies any entry of $\{Q^k_{\text{m}}, Q^k_{\text{s}}\}$.
This is a deliberate design choice: a previous implementation cached unseen-entry
fallbacks into each $Q^k$ to avoid recomputing the heuristic, but this caused
the dictionaries to grow at rate $O(\text{unique states per battle} \times K)$,
exceeding memory in a 100{,}000-battle evaluation by the 30{,}000-battle mark.
Our read-only formulation keeps the inference memory footprint at $O(1)$ per
battle and admits the longer evaluations our experiments require.

\subsection{Training the ensemble members}
\label{sec:method_training}

% K = 30, seeded by base_seed + k
% Hyperparameters: alpha=0.1, gamma=0.99, lambda=0.7, eps=0.05 fixed
% 20-Pokemon pool, gen4ou format, IndexedTeambuilder (deterministic teams per battle)
% 1M battles per agent, parallelism

Each of $K = 30$ members is trained independently with seed $\text{seed}_k = s_0 + k$,
$\alpha = 0.1$, $\gamma = 0.99$, $\lambda = 0.7$, and a fixed exploration rate
$\varepsilon = 0.05$. Members observe and learn from a fixed 20-Pokemon OU pool from
which the IndexedTeambuilder constructs deterministic teams indexed by battle number
(so battle $N$ produces the same matchup across members). Each member trains for
$1{,}000{,}000$ battles, for a total of $30{,}000{,}000$ training battles across the
ensemble. Training is parallelized across eight AWS EC2 c5.4xlarge instances over
fifteen days.

% Reproducibility note: the entire pipeline is on GitHub.
The training pipeline and trained checkpoints are released at
\url{https://github.com/arzaansingh/poke_rl}.
